\SourceLeaf{266}
\fancyhead[L]{\small UNTERSUCHUNGEN ÜBER GEGENSTÄNDE DER HÖHERN GEODAESIE}
\renewcommand{\Article}[1]{\par\medskip\Needspace{4\baselineskip}\pdfbookmark[1]{Art. #1}{geodaesie1-art-#1}\centerline{\textbf{#1.}}\nobreak\smallskip}
\begin{center}
UNTERSUCHUNGEN\par\bigskip
ÜBER\par\bigskip
{\large GEGENSTÄNDE DER HÖHERN GEODAESIE.}\par\medskip
\rule{.65in}{.3pt}
\end{center}
\noindent\hspace*{\parindent}Bei den, zum Theil von mir selbst, zum Theil unter meiner Leitung, ausgeführten über das ganze Königreich Hannover sich erstreckenden trigonometrischen Vermessungen sind, sowohl in Beziehung auf die Art, wie die Messungen angestellt wurden, als noch mehr in Beziehung auf ihre nachherige mathematische Behandlung und ihre Verarbeitung zu Resultaten, Wege eingeschlagen, die von den sonst gewöhnlichen abweichen. Mein früher gehegter Vorsatz, nach völliger Beendigung der Messungen diese nebst allen von mir angewandten Verfahrungsarten in einem besondern Werke darzulegen, hat, aus Ursachen, deren Auseinandersetzung nicht hieher gehört, bisher nicht zur Ausführung kommen können, und ich wähle daher das Auskunftsmittel, das im theoretischen Theile mir eigenthümliche in einer Reihe von Abhandlungen bekannt zu machen, um so lieber, weil ich auf diese Weise die Freiheit behalte, mit Ausführlichkeit manche Untersuchungen zu entwickeln, welche ein selbstständiges Interesse darbieten und mit den übrigen in enger Verwandtschaft stehen, auch wenn von denselben bei meinen Messungen keine unmittelbare Anwendung gemacht ist. Dies gilt namentlich von dem grössten Theile des Inhalts der gegenwärtigen ersten Abhandlung.
